% Generated by publication/build_sections.py; do not edit by hand.
\paragraph{Informal verdict.}
\textbf{A --- Complete solution}, with high confidence.

\paragraph{Lean coverage.}
Lean proves the complete necessary-and-sufficient Q730 classification, including prescribed diagonal order, rank-deficient prefix inequalities, and the full-rank determinant equality.

\subsection{Audited informal answer}
\subsubsection{1. Faithful restatement}

Let \(M\in M_n(\mathbb C)\) have rank \(r\), and let \(t_1,\dots,t_r>0\). Determine when there exist unitary matrices \(U,V\in U(n)\) such that

\[
UMV=
\begin{pmatrix}
T&B\\
0&0
\end{pmatrix},
\]

where \(T\in M_r(\mathbb C)\) is upper triangular with diagonal

\[
\operatorname{diag}(T)=(t_1,\dots,t_r),
\]

and \(B\in M_{r,n-r}(\mathbb C)\). The lower-left and lower-right blocks are zero matrices of sizes \((n-r)\times r\) and \((n-r)\times(n-r)\), respectively.

The PDF adds that an earlier answer treated the full-rank case \(|\det M|=1\) with \(t_1=\cdots=t_n=1\). 

The phrase ``Soit des réels \(t_1,\ldots,t_r\)'' is a minor grammatical typo; mathematically it clearly means ``Let \(t_1,\ldots,t_r\) be positive real numbers.'' The use of \(UMV\), rather than the more usual \(UMV^*\), causes no issue because \(V\mapsto V^*\) is a bijection of the unitary group.

\medskip\hrule\medskip

\subsubsection{2. Final classification}

Let

\[
\sigma_1\ge \sigma_2\ge\cdots\ge \sigma_r>0
\]

be the nonzero singular values of \(M\). Let

\[
\tau_1\ge\tau_2\ge\cdots\ge\tau_r>0
\]

be the decreasing rearrangement of \(t_1,\dots,t_r\).

\paragraph{Theorem}

The required unitary matrices (U,V) exist precisely under the following conditions.

\paragraph{Rank-deficient case \(r<n\)}

They exist if and only if

\[
\boxed{\quad
\prod_{i=1}^k \tau_i
\le
\prod_{i=1}^k \sigma_i
\qquad (1\le k\le r).
\quad}
\tag{1}
\]

Thus \((\tau_i)\) is weakly log-majorized by \((\sigma_i)\).

\paragraph{Full-rank case \(r=n\)}

They exist if and only if

\[
\boxed{\quad
\prod_{i=1}^k \tau_i
\le
\prod_{i=1}^k \sigma_i
\qquad (1\le k<n),
\quad}
\tag{2}
\]

and

\[
\boxed{\quad
\prod_{i=1}^n t_i
=
%
= \prod_{i=1}^n \sigma_i
%
|\det M|.
\quad}
\tag{3}
\]

Thus \((\tau_i)\) is log-majorized by \((\sigma_i)\).

The order \(t_1,\dots,t_r\) requested on the diagonal imposes no additional condition: once the eigenvalues have been realized, the unitary Schur decomposition can put them on the diagonal in any prescribed order.

The central eigenvalue--singular-value criterion is traditionally called the Weyl--Horn theorem. A complete proof sufficient for the present problem is included below rather than used as a black box. ([Journals de l'Université du Wyoming][1])

\medskip\hrule\medskip

\subsection{3. Preparatory results}

\subsubsection{3.1. Two-sided unitary equivalence}

\paragraph{Lemma 1}

Two matrices \(A,B\in M_n(\mathbb C)\) have the same singular values, counted with multiplicity, if and only if there exist unitary matrices (P,Q) such that

\[
B=PAQ.
\]

\paragraph{Proof}

The forward implication follows from the singular value decompositions

\[
A=P_A\Sigma Q_A^*,
\qquad
B=P_B\Sigma Q_B^*,
\]

with the same nonnegative diagonal matrix \(\Sigma\). Then

\[
B=(P_BP_A^*),A,(Q_AQ_B^*).
\]

Both factors are unitary.

Conversely, if \(B=PAQ\) with (P,Q) unitary, then

\[
B^*B=Q^*A^*AQ,
\]

so \(A^*A\) and \(B^*B\) have the same eigenvalues. Hence \(A\) and \(B\) have the same singular values. \(\square\)

\medskip\hrule\medskip

\subsubsection{3.2. Ordered unitary triangularization}

\paragraph{Lemma 2}

Let \(A\in M_n(\mathbb C)\), and let \(\lambda_1,\dots,\lambda_n\) be its eigenvalues, listed with algebraic multiplicity in any prescribed order. There exists a unitary \(W\) such that \(W^*AW\) is upper triangular with diagonal

\[
(\lambda_1,\dots,\lambda_n).
\]

\paragraph{Proof}

Choose a unit eigenvector \(x_1\) for \(\lambda_1\) and extend it to an orthonormal basis. In that basis,

\[
A=
\begin{pmatrix}
\lambda_1&*\\
0&A_1
\end{pmatrix}.
\]

The characteristic polynomial factors as

\[
\chi_A(z)=(z-\lambda_1)\chi_{A_1}(z),
\]

so the eigenvalues of \(A_1\) are \(\lambda_2,\dots,\lambda_n\), with multiplicity. Apply the same argument recursively to \(A_1\). \(\square\)

\medskip\hrule\medskip

\subsection{4. Eigenvalues versus singular values}

We now prove the exact form of the Weyl--Horn theorem needed here.

\subsubsection{Theorem 3 --- Weyl--Horn criterion}

Let

\[
s_1\ge\cdots\ge s_n\ge0
\]

and let \(\lambda_1,\dots,\lambda_n\in\mathbb C\), ordered so that

\[
|\lambda_1|\ge\cdots\ge|\lambda_n|.
\]

There exists a matrix \(A\in M_n(\mathbb C)\) whose singular values are \(s_1,\dots,s_n\) and whose eigenvalues are \(\lambda_1,\dots,\lambda_n\) if and only if

\[
\prod_{i=1}^k|\lambda_i|
\le
\prod_{i=1}^k s_i
\qquad (1\le k<n),
\tag{4}
\]

and

\[
\prod_{i=1}^n|\lambda_i|
=
%
\prod_{i=1}^n s_i.
\tag{5}
\]

\medskip\hrule\medskip

\subsubsection{Proof: necessity}

For \(1\le k\le n\), consider the induced operator

\[
\bigwedge\nolimits^k A
\]

on \(\bigwedge^k\mathbb C^n\).

By upper triangularizing \(A\), one sees that the eigenvalues of \(\bigwedge^k A\) are all products

\[
\lambda_{i_1}\cdots\lambda_{i_k},
\qquad
1\le i_1<\cdots<i_k\le n.
\]

Because the moduli of the \(\lambda_i\) are decreasing,

\[
\rho\!\left(\bigwedge\nolimits^k A\right)
=
%
\prod_{i=1}^k|\lambda_i|,
\]

where \(\rho\) denotes the spectral radius.

On the other hand, applying the exterior-power functor to a singular value decomposition of \(A\) shows that the singular values of \(\bigwedge^k A\) are the products

\[
s_{i_1}\cdots s_{i_k}.
\]

Consequently,

\[
\left\lvert \bigwedge\nolimits^k A\right\rvert 
=
%
\prod_{i=1}^k s_i.
\]

Since the spectral radius does not exceed the operator norm,

\[
\prod_{i=1}^k|\lambda_i|
\le
\prod_{i=1}^k s_i.
\]

This proves \(4\). Finally,

\[
\prod_{i=1}^n|\lambda_i|
=
%
= |\det A|
%
= \sqrt{\det(A^*A)}
%
\prod_{i=1}^n s_i,
\]

which is \(5\).

\medskip\hrule\medskip

\subsubsection{Proof: sufficiency}

We argue by induction on \(n\).

For \(n=1\), condition \(5\) says \(s_1=|\lambda_1|\), so the matrix \((\lambda_1)\) works.

Assume the result known in dimension \(n-1\).

\paragraph{Case 1: \(\lambda_1=0\)}

Then all the \(\lambda_i\) vanish. Equation \(5\) implies

\[
s_n=0.
\]

Consider the weighted shift

\[
A=
\begin{pmatrix}
0&s_1&0&\cdots&0\\
0&0&s_2&\ddots&\vdots\\
\vdots&\ddots&\ddots&\ddots&0\\
0&\cdots&0&0&s_{n-1}\\
0&\cdots&\cdots&0&0
\end{pmatrix}.
\]

It is nilpotent, so all its eigenvalues are zero. Moreover,

\[
A^*A=\operatorname{diag}(0,s_1^2,\dots,s_{n-1}^2),
\]

so its singular values are

\[
s_1,\dots,s_{n-1},0=s_n.
\]

\paragraph{Case 2: \(\lambda_1\neq0\)}

Put

\[
a=|\lambda_1|>0.
\]

Condition \(4\) for \(k=1\) gives \(a\le s_1\). We also have \(a\ge s_n\): otherwise

\[
|\lambda_i|\le a<s_n\le s_i
\]

for every \(i\), contradicting the product equality \(5\).

Hence there is an index \(j\in{2,\dots,n}\) such that

\[
s_{j-1}\ge a\ge s_j.
\tag{6}
\]

Define

\[
d=\frac{s_1s_j}{a}.
\tag{7}
\]

Consider

\[
H=
\begin{pmatrix}
\lambda_1&\mu\\
0&d
\end{pmatrix},
\]

where

\[
\mu^2
=
%
s_1^2+s_j^2-a^2-d^2.
\]

Using \(7\),

\[
\begin{aligned}
\mu^2
&=
s_1^2+s_j^2-a^2-\frac{s_1^2s_j^2}{a^2}\\
&=
\frac{(s_1^2-a^2)(a^2-s_j^2)}{a^2}
\ge0
\end{aligned}
\]

by \(6\). We take \(\mu\ge0\).

Now

\[
\operatorname{tr}(H^*H)
=
%
= a^2+d^2+\mu^2
%
s_1^2+s_j^2
\]

and

\[
\det(H^*H)=|\det H|^2
=a^2d^2=s_1^2s_j^2.
\]

Thus the two eigenvalues of \(H^*H\) are \(s_1^2,s_j^2\), and the singular values of \(H\) are \(s_1,s_j\).

Consider the following multiset of \(n-1\) nonnegative numbers:

\[
\mathcal S'
=
%
\left\{
d,s_2,\dots,s_{j-1},s_{j+1},\dots,s_n
\right\}.
\]

Let

\[
y_1\ge\cdots\ge y_{n-1}
\]

be its decreasing rearrangement. We claim that the singular-value data \(y_1,\dots,y_{n-1}\) and eigenvalue data \(\lambda_2,\dots,\lambda_n\) satisfy the induction hypotheses.

First,

\[
\begin{aligned}
\prod_{i=2}^n|\lambda_i|
&=
\frac{1}{a}\prod_{i=1}^n s_i\\
&=
d\prod_{\substack{2\le i\le n\\i\ne j}}s_i\\
&=
\prod_{\ell=1}^{n-1}y_\ell.
\end{aligned}
\tag{8}
\]

It remains to prove the partial-product inequalities.

Let \(1\le m\le j-2\). Since \(s_2,\dots,s_{m+1}\) are members of \(\mathcal S'\), and each of them is at least \(a\), we have

\[
\begin{aligned}
\prod_{i=2}^{m+1}|\lambda_i|
&\le a^m\\
&\le \prod_{i=2}^{m+1}s_i\\
&\le \prod_{\ell=1}^m y_\ell.
\end{aligned}
\tag{9}
\]

Now let \(j-1\le m\le n-2\). Applying \(4\) with \(k=m+1\) gives

\[
a\prod_{i=2}^{m+1}|\lambda_i|
\le
\prod_{i=1}^{m+1}s_i.
\]

Therefore

\[
\begin{aligned}
\prod_{i=2}^{m+1}|\lambda_i|
&\le
\frac1a\prod_{i=1}^{m+1}s_i\\
&=
d\prod_{\substack{2\le i\le m+1\\i\ne j}}s_i\\
&\le
\prod_{\ell=1}^m y_\ell.
\end{aligned}
\tag{10}
\]

The final inequality holds because the preceding expression is the product of \(m\) members of \(\mathcal S'\), while \(y_1\cdots y_m\) is the product of its \(m\) largest members.

By the induction hypothesis, there exists \(C\in M_{n-1}(\mathbb C)\) with singular values \(y_1,\dots,y_{n-1}\) and eigenvalues \(\lambda_2,\dots,\lambda_n\).

Let

\[
D=\operatorname{diag}
\bigl(
d,s_2,\dots,s_{j-1},s_{j+1},\dots,s_n
\bigr).
\]

Since \(C\) and \(D\) have the same singular values, Lemma 1 gives unitary \(P,Q\in U(n-1)\) such that

\[
C=PDQ.
\]

Define

\[
A_0=
\begin{pmatrix}
\lambda_1&\mu e_1^*\\
0&D
\end{pmatrix},
\]

where \(e_1\) is the first standard basis vector of \(\mathbb C^{n-1}\). The matrix \(A_0\) is the direct sum, up to the displayed ordering of coordinates, of the \(2\times2\) matrix \(H\) and the diagonal entries \(s_i\) for \(i\ne1,j\). Hence its singular values are \(s_1,\dots,s_n\).

Finally set

\[
A=(1\oplus P)A_0(1\oplus Q).
\]

Its singular values remain \(s_1,\dots,s_n\), and

\[
A=
\begin{pmatrix}
\lambda_1&\mu e_1^*Q\\
0&C
\end{pmatrix}.
\]

Thus its eigenvalues are \(\lambda_1\) together with the eigenvalues \(\lambda_2,\dots,\lambda_n\) of \(C\). The induction is complete. \(\square\)

\medskip\hrule\medskip


\subsubsection{5.1. Necessity}

Suppose that the requested unitary matrices exist, and set

\[
A=UMV=
\begin{pmatrix}
T&B\\
0&0
\end{pmatrix}.
\]

Because multiplication on the left and right by unitaries preserves singular values, the singular values of \(A\) are

\[
\sigma_1,\dots,\sigma_r,
\underbrace{0,\dots,0}_{n-r}.
\]

Since \(A\) is block upper triangular and \(T\) is upper triangular with diagonal \(t_1,\dots,t_r\), the eigenvalues of \(A\) are

\[
t_1,\dots,t_r,
\underbrace{0,\dots,0}_{n-r}.
\]

Order the positive eigenvalues decreasingly as \(\tau_1\ge\cdots\ge\tau_r\). The necessary part of Theorem 3 gives

\[
\prod_{i=1}^k\tau_i
\le
\prod_{i=1}^k\sigma_i,
\qquad 1\le k\le r.
\]

If \(r=n\), the determinant identity additionally gives

\[
\prod_{i=1}^n t_i
=
%
= |\det A|
%
= |\det M|
%
\prod_{i=1}^n\sigma_i.
\]

This proves necessity.

\medskip\hrule\medskip

\subsubsection{5.2. Sufficiency}

Assume first that \(r>0\), and define the full lists

\[
s=
(\sigma_1,\dots,\sigma_r,0,\dots,0)
\]

and

\[
\lambda=
(\tau_1,\dots,\tau_r,0,\dots,0).
\]

If \(r<n\), the assumed inequalities give the Weyl--Horn conditions for \(k\le r\). For every \(k\ge r+1\),

\[
\prod_{i=1}^k|\lambda_i|
=
%
= 0
%
\prod_{i=1}^k s_i.
\]

In particular, the full products agree.

If \(r=n\), the assumed partial-product inequalities and the determinant equality are exactly the Weyl--Horn conditions.

Theorem 3 therefore provides a matrix \(A_0\in M_n(\mathbb C)\) with the same singular values as \(M\) and with eigenvalues

\[
t_1,\dots,t_r,0,\dots,0
\]

as a multiset.

By Lemma 1, there exist unitary matrices \(U_0,V_0\) such that

\[
A_0=U_0MV_0.
\tag{11}
\]

Use Lemma 2 to choose a unitary \(W\) such that

\[
R=W^*A_0W
\]

is upper triangular with diagonal in the particular order

\[
t_1,\dots,t_r,0,\dots,0.
\]

Partition \(R\) as

\[
R=
\begin{pmatrix}
T&B\\
0&C
\end{pmatrix},
\]

where \(T\in M_r(\mathbb C)\). Then \(T\) is upper triangular with diagonal \(t_1,\dots,t_r\), hence \(T\) is invertible.

The rank of \(R\) equals the rank of \(M\), namely \(r\). But right multiplication by the invertible matrix

\[
\begin{pmatrix}
I_r&-T^{-1}B\\
0&I_{n-r}
\end{pmatrix}
\]

gives

\[
R
\begin{pmatrix}
I_r&-T^{-1}B\\
0&I_{n-r}
\end{pmatrix}
=
%
\begin{pmatrix}
T&0\\
0&C
\end{pmatrix}.
\]

Therefore

\[
\operatorname{rank}R
=
%
= \operatorname{rank}T+\operatorname{rank}C
%
r+\operatorname{rank}C.
\]

Since \(\operatorname{rank}R=r\), necessarily \(C=0\).

Combining this with \(11\),

\[
R
=
%
(W^*U_0),M,(V_0W).
\]

Both \(W^*U_0\) and \(V_0W\) are unitary, and \(R\) has exactly the required form.

If \(r=0\), then \(M=0\), the list \(t_1,\dots,t_r\) is empty, and the assertion is trivially true.

This completes the proof. \(\square\)

\medskip\hrule\medskip

\subsection{6. Boundary cases and checks}

\subsubsection{Rank one}

Suppose \(r=1<n\), and let \(\sigma_1\) be the unique nonzero singular value. The condition reduces to

\[
t_1\le \sigma_1.
\]

This is visibly sufficient: the matrix

\[
\begin{pmatrix}
t_1&\sqrt{\sigma_1^2-t_1^2}&0&\cdots&0\\
0&0&0&\cdots&0\\
\vdots&\vdots&\vdots&&\vdots
\end{pmatrix}
\]

has singular values \(\sigma_1,0,\dots,0\) and the required form.

For \(n=r=1\), there is no \(B\)-block, and the condition becomes the equality

\[
t_1=\sigma_1=|M|.
\]

This shows explicitly why the product equality disappears as soon as a zero singular value is available.

\subsubsection{Higher partial products are genuinely necessary}

It is not enough to require only \(t_i\le\sigma_1\). For example, suppose the nonzero singular values are

\[
(\sigma_1,\sigma_2)=(3,1)
\]

and the desired diagonal is ((2,2)). Each desired diagonal entry is at most \(3\), but

\[
2\cdot2=4>3\cdot1.
\]

Hence no such unitary reduction exists.

\subsubsection{In full rank, the determinant condition alone is not enough}

Take singular values

\[
(4,\tfrac14)
\]

and desired diagonal

\[
(5,\tfrac15).
\]

Both products equal \(1\), but

\[
5>4,
\]

so the first partial-product inequality fails. Thus determinant equality must be supplemented by all the inequalities \(2\).

\subsubsection{Recovery of the earlier special case}

Suppose \(r=n\), \(|\det M|=1\), and \(t_1=\cdots=t_n=1\). Then

\[
\prod_{i=1}^n\sigma_i=1.
\]

For every \(k<n\),

\[
\prod_{i=1}^k\sigma_i\ge1.
\]

Indeed, if this product were \(<1\), then \(\sigma_k<1\), so every later singular value would also be \(<1\), forcing the total product to be \(<1\), a contradiction. Consequently,

\[
1=\prod_{i=1}^k t_i
\le
\prod_{i=1}^k\sigma_i,
\]

and equality holds for \(k=n\). The theorem therefore recovers exactly the affirmative case mentioned in the PDF.

\medskip\hrule\medskip

\subsubsection{Final answer}

Let \(\sigma_1\ge\cdots\ge\sigma_r>0\) be the nonzero singular values of \(M\), and let \(\tau_1\ge\cdots\ge\tau_r\) be the decreasing rearrangement of \(t_1,\dots,t_r\).

\[
\boxed{
\begin{array}{ll}
r<n:
&
\displaystyle
\prod_{i=1}^k\tau_i
\le
\prod_{i=1}^k\sigma_i
\quad\text{for every }1\le k\le r;
\\[1.4em]
r=n:
&
\displaystyle
\prod_{i=1}^k\tau_i
\le
\prod_{i=1}^k\sigma_i
\quad\text{for }1\le k<n,
\\[0.8em]
&
\displaystyle
\prod_{i=1}^n t_i=|\det M|.
\end{array}
}
\]

These conditions are both necessary and sufficient, and the diagonal entries can be placed in the exact prescribed order \(t_1,\dots,t_r\).

[1]: https://journals.uwyo.edu/index.php/ela/article/download/823/823/823 "A. Horn's result on matrices with prescribed singular values and eigenvalues"
